\section*{Bab \ref{ch:g12:climate-change} --- Perubahan Iklim: Pengamatan, Penyebab, dan Proyeksi}

\begin{solution}{exo:g12:climate-change:1}
Perubahan yang ditimbulkannya pada aliran bersih energi yang memasuki sistem
iklim di puncak atmosfer, sebelum suhu permukaan menanggapinya; diukur dalam
\unit{W/m^2}. Positif: planet memperoleh energi dan menghangat; negatif:
planet kehilangan energi dan mendingin.
\end{solution}

\begin{solution}{exo:g12:climate-change:2}
Paling tinggi pada bulan Mei, paling rendah pada bulan September. Pada musim
semi dan musim panas, vegetasi Belahan Bumi Utara, tempat sebagian besar
tumbuhan darat tumbuh, menyerap lebih banyak karbon dioksida melalui
fotosintesis daripada yang dilepaskan tumbuhan dan tanah; pada musim gugur dan
musim dingin, pembusukan dan respirasi yang menang, dan konsentrasinya naik
lagi.
\end{solution}

\begin{solution}{exo:g12:climate-change:3}
\emph{Sangat mungkin}: peluang 90 sampai 100\,\%; \emph{mungkin}: 66 sampai
100\,\%. $[0.95 \text{ sampai } 1.20]$ adalah rentang sangat mungkin: ada peluang 90\,\% bahwa
nilai sebenarnya terletak antara 0.95 dan 1.20.
\end{solution}

\begin{solution}{exo:g12:climate-change:4}
91\,\%. Air memuai ketika menghangat: samudra yang lebih hangat menempati
volume yang lebih besar, sehingga permukaannya naik walaupun tidak ada air
yang ditambahkan (es darat yang mencair menambahkan air di atas kenaikan ini).
\end{solution}

\begin{solution}{exo:g12:climate-change:5}
Uap air: udara yang lebih hangat menampung lebih banyak uap air, sebuah gas
rumah kaca, yang menambah pemanasan. Salju dan es: pencairannya menyingkapkan
permukaan yang lebih gelap, yang menyerap lebih banyak sinar Matahari.
Tanggapan negatifnya: planet yang lebih hangat memancarkan lebih banyak
inframerah ke angkasa, yang mengembalikannya ke dalam keseimbangan.
\end{solution}

\begin{solution}{exo:g12:climate-change:6}
$427.35 - 315.98 = \qty{111.37}{ppm}$; $111.37/315.98 \approx 0.352$, sebuah
kenaikan sekitar 35\,\%. Selama 66 tahun: $111.37/66 \approx
\qty{1.7}{ppm}$ per tahun secara rata-rata.
\end{solution}

\begin{solution}{exo:g12:climate-change:7}
$F = 3.93 \times \ln(427.35/278)/\ln 2 = 3.93 \times 0.4300/0.6931 \approx
\qty{2.44}{W/m^2}$.
\end{solution}

\begin{solution}{exo:g12:climate-change:8}
$3.7 \times 30 = \qty{111}{mm}$, sekitar \qty{11}{cm}. $3.7/1.3 \approx 2.8$:
hampir tiga kali lebih cepat.
\end{solution}

\begin{solution}{exo:g12:climate-change:9}
$\Delta T = 2.72/1.16 \approx \qty{2.3}{\degreeCelsius}$. Iklim belum berada
dalam kesetimbangan: samudra, yang menyerap sekitar 90\,\% energi tambahan
itu, menghangat dengan lambat, sehingga permukaan belum menyusul \omterm{def:g12:climate-change:forcing}{gayanya}
(sistem itu masih memperoleh sekitar \qty{0.8}{W/m^2}).
\end{solution}

\begin{solution}{exo:g12:climate-change:10}
Matahari dan gunung api paling banyak menjelaskan \qty{0.1}{\degreeCelsius}
dari pemanasan \qty{1.06}{\degreeCelsius}: keduanya tidak mungkin menjadi
penyebabnya. Penyebab alami yang lain adalah variabilitas internal iklim
(fluktuasi alaminya, seperti pertukaran panas antara samudra dan atmosfer);
penilaian menempatkan sumbangannya pada $-0.2$ sampai $+\qty{0.2}{\degreeCelsius}$,
juga jauh terlalu kecil.
\end{solution}

\begin{solution}{exo:g12:climate-change:11}
Sangat mungkin di bawah \qty{2}{\degreeCelsius}: hanya SSP1-1.9 (rentang
sangat mungkin 1.0 sampai 1.8). Sangat mungkin di atasnya: SSP2-4.5 (2.1
sampai 3.5), SSP3-7.0 (2.8 sampai 4.6), dan SSP5-8.5 (3.3 sampai 5.7). Untuk
SSP1-2.6 (1.3 sampai 2.4) tidak satu pun dari keduanya sangat mungkin.
\end{solution}

\begin{solution}{exo:g12:climate-change:12}
\omterm{def:g12:climate-change:forcing}{Gaya} aerosol adalah $-0.22 - 0.84 = \qty{-1.06}{W/m^2}$; tanpa \omterm{def:g12:climate-change:forcing}{gaya} itu
totalnya menjadi $2.72 + 1.06 = \qty{3.78}{W/m^2}$. Pemanasan kesetimbangan
bertambah $1.06/1.16 \approx \qty{0.9}{\degreeCelsius}$.
\end{solution}

\begin{solution}{exo:g12:climate-change:13}
Karena $\ln(2C) - \ln C = \ln 2$ berapa pun $C$,
\[
  F(556) - F(278) = \frac{3.93\,\ln 2}{\ln 2} = \qty{3.93}{W/m^2}, \qquad
  F(1112) - F(556) = \frac{3.93\,\ln 2}{\ln 2} = \qty{3.93}{W/m^2}:
\]
\omterm{def:g12:climate-change:forcing}{gaya} yang sama untuk setiap penggandaan. Total pada \qty{1112}{ppm}: $2 \times 3.93 =
\qty{7.86}{W/m^2}$.
\end{solution}

\begin{solution}{exo:g12:climate-change:14}
$\dfrac{\dd F}{\dd C} = \dfrac{F_{2\times}}{C \ln 2}$. Pada \qty{280}{ppm}:
$3.93/(280 \times 0.6931) \approx \qty{0.0202}{W/m^2}$ per ppm; pada
\qty{420}{ppm}: $\approx \qty{0.0135}{W/m^2}$ per ppm. Setiap ppm tambahan
menambahkan \omterm{def:g12:climate-change:forcing}{gaya} yang makin kecil ketika konsentrasinya bertambah (1.5 kali
lebih kecil pada 420 daripada pada 280), tetapi konsentrasinya terus naik,
sehingga \omterm{def:g12:climate-change:forcing}{gayanya} tetap bertambah.
\end{solution}

\begin{solution}{exo:g12:climate-change:15}
$3.93/3.22 \approx \qty{1.2}{\degreeCelsius}$ tanpa \omterm{def:g12:climate-change:feedback}{umpan balik};
$3.93/1.16 \approx \qty{3.4}{\degreeCelsius}$ dengan \omterm{def:g12:climate-change:feedback}{umpan balik}. Faktor
penguatannya: $3.22/1.16 \approx 2.8$.
\end{solution}

\begin{solution}{pb:g12:climate-change:1}
\textbf{1.} $3.37 + 3.51 = \qty{6.88}{ppm}$, sekitar \qty{7}{ppm}.

\textbf{2.} Dari musim semi sampai akhir musim panas, tumbuhan Belahan Bumi
Utara menarik karbon dioksida dari udara melalui fotosintesis lebih cepat
daripada respirasi dan pembusukan mengembalikannya: konsentrasinya turun
sampai bulan September. Pada musim gugur dan musim dingin neracanya berbalik,
dan konsentrasinya naik kembali ke maksimumnya pada bulan Mei.

\textbf{3.} Belahan Bumi Selatan memiliki daratan yang jauh lebih sedikit, dan
dengan demikian vegetasi yang jauh lebih sedikit, daripada Belahan Bumi Utara;
sebagian besar permukaannya adalah samudra.

\textbf{4.} Karena tinggi, observatorium itu terletak di atas lapisan udara
tercemar di dekat tanah; jauh dari benua dan kota, udara yang dicuplikinya
tercampur baik dan mewakili seluruh atmosfer, bukan sumber setempat.

\textbf{5.} $(325.68 - 316.91)/10 \approx \qty{0.88}{ppm}$ per tahun.

\textbf{6.} $(427.35 - 401.01)/10 \approx \qty{2.63}{ppm}$ per tahun.

\textbf{7.} $2.634/0.877 \approx 3.0$: laju pertambahannya menjadi tiga kali
lipat.

\textbf{8.} $(427.35 - 278)/278 \approx 0.537$: sekitar 54\,\%.

\textbf{9.} $427.35 - 300 \approx \qty{127}{ppm}$.

\textbf{10.} $F = 3.93 \times \ln(427.35/278)/\ln 2 \approx
\qty{2.44}{W/m^2}$.

\textbf{11.} $F(316.91) = 3.93 \times \ln(316.91/278)/\ln 2 \approx
\qty{0.74}{W/m^2}$; kenaikannya $2.44 - 0.74 \approx \qty{1.70}{W/m^2}$.

\textbf{12.} Satu \unit{W/m^2} lagi mengalikan konsentrasinya dengan
$2^{1/3.93} \approx 1.193$: $427.35 \times 1.193 \approx \qty{510}{ppm}$.

\textbf{13.} $(509.8 - 427.35)/2.634 \approx 31$ tahun.

\textbf{14.} $2 \times 278 = \qty{556}{ppm}$; \omterm{def:g12:climate-change:forcing}{gayanya} $F_{2\times} =
\qty{3.93}{W/m^2}$.

\textbf{15.} Menurut definisi sensitivitas itu, \qty{3}{\degreeCelsius}
(mungkin 2.5 sampai \qty{4}{\degreeCelsius}).

\textbf{16.} $\lambda = 3.93/3 \approx \qty{1.31}{W.m^{-2}.\degreeCelsius^{-1}}$.
Jumlah \omterm{def:g12:climate-change:feedback}{umpan balik-umpan baliknya}, \qty{1.16}{W.m^{-2}.\degreeCelsius^{-1}}, memberikan
$3.93/1.16 \approx \qty{3.4}{\degreeCelsius}$, sedikit lebih besar; perkiraan
terbaik penilaian menggabungkan fisika ini dengan pemanasan yang teramati dan
iklim masa lalu, dan kedua nilai itu terletak di dalam rentang mungkin.

\textbf{17.} Samudra menyimpan sebagian besar energi tambahan itu dan
menghangat dengan lambat: permukaan baru mencapai pemanasan kesetimbangannya
setelah beberapa dekade sampai beberapa abad.

\textbf{18.} Perhitungan itu mengandaikan satu cerita tentang emisi masa depan
(pertambahan sepuluh tahun terakhir dipertahankan selamanya): itulah sebuah
\omterm{def:g12:climate-change:scenario}{skenario}. Sebuah \omterm{def:g12:climate-change:scenario}{proyeksi} menghitung tanggapan iklim dan siklus karbon
terhadap suatu \omterm{def:g12:climate-change:scenario}{skenario emisi} dengan model-model fisika, dan memberikan
rentang serta kemungkinan.

\textbf{19.} $(556 - 427.35)/2.634 \approx 48.8$ tahun setelah 2025: sekitar
\textbf{2074}.
\end{solution}
